\documentclass[12pt]{article}
\input{tcilatex}\textheight 22.5cm \textwidth 16cm\setlength{\oddsidemargin}{-0.1cm} \setlength{\topmargin}{-1.5cm}
\begin{document}

\noindent Lars Ljungqvist  \hfill  Macroeconomic Theory I
 
\noindent New York University                               \hskip7.55cm Fall 2023



\vskip1cm


\centerline{{\bf Final Exam} (100 points)}
\vskip.6cm

\vskip.5cm
\noindent
{\it Explain your answers both with mathematics and in words. 
If you find any ambiguity in the questions, state your clarifying 
assumptions and proceed with the analysis.}


\vskip1cm

\noindent
Consider an infinite-horizon pure endowment economy.  
In each period $t\geq 0$, there is a realization of 
a stochastic event $s_{t}\in \mathbf{S}$.
The history of events up and until time $t$
is denoted $s^t = [s_0, s_1, \ldots, s_t]$. The 
unconditional probability of observing a particular 
sequence of events $s^t$ is given by a probability 
measure $\pi_t(s^t)$. The economy starts at time $0$ 
with a known $s_0$, i.e., $\pi_0(s_0)=1$ for that
particular $s_0$. 

There are equal numbers of two types of consumers, $i=1,2$. 
Consumers of type $i$ order consumption
streams of the one good according to
$$ \sum_{t=0}^\infty \sum_{s^t} \beta^t \pi_t(s^t)
\left\{\alpha c_{it}(s^t) 
- \frac{[c_{it}(s^t)]^2}{2}\right\}, \hskip.6cm 
 \beta\in(0,1),\;\, \alpha>1, $$
where $c_{it}(s^t)\geq 0$ is the consumption of a type $i$ consumer
after history $s^t$. The good is tradable but nonstorable. 
At time $t$ and history $s^t$, consumer $1$ receives
the endowment $y_{1t}(s^t) \in (0,1)$ and consumer $2$
receives $y_{2t}(s^t) = 1 - y_{1t}(s^t)$.


\begin{itemize}

\item[a.] \textbf{(8 points)} 
Solve the social planner problem
with Pareto weights $\lambda_1$ and $\lambda_2$ for
consumers of type 1 and 2, respectively. Find the 
range of relative Pareto weights, $\lambda_1/\lambda_2$, that
fully maps out the Pareto frontier.
\medskip

\item[b.] \textbf{(4 points)} 
Define a competitive equilibrium with 
time $0$ trading. 
\medskip

\item[c.] \textbf{(20 points)} 
Compute a time $0$ trading  
equilibrium, i.e., find an allocation 
$\{c_{it}(s^t); \forall i,$ $\forall t, \forall s^t\}$ 
and prices $\{q^0_t(s^t); \forall t, \forall s^t\}$
expressed in terms of primitives.
\medskip

\item[d.] \textbf{(3 points)} 
Define a competitive equilibrium with 
sequential trading. 


\end{itemize}

\vskip.5cm

\noindent
{\bf From hereon}, we assume that the stochastic event $s_t$
can take on two values, $s_{t}\in \{0,1\}$, with
the following probabilities. 
The economy starts with $s_0=0$, i.e., $\pi_0(0)=1$. Next
period, either $s_1=0$ or $s_1=1$ is realized with
probabilities $\bar \pi \in(0,1)$ and $1-\bar \pi$,
respectively: $\pi_1(0, 0)=\bar \pi$ and
$\pi_1(0, 1)= 1 - \bar \pi$. Thereafter, the realization
of the stochastic event in period~1 persists forever, 
$s_t=s_1$ for all $t>1$. Thus, for $t\geq 1$, there are
just two possible paths of $s^t$, namely
$s^t = [0, 0, 0, \ldots 0, 0]$ and
$s^t = [0, 1, 1, \ldots 1, 1]$. The endowment process 
of a consumer of type~1 is given by
\begin{equation}
y_{1t}(s^t) = \left\{ \begin{array}{cc}
    \bar y_{10}\in (0,1),     & \hbox{\rm if  } s_t=0;  \\
    \bar y_{11}\in (0,1),     & \hbox{\rm if  } s_t=1;
\end{array} \right.
\end{equation}
and as stated earlier, the endowment of a consumer of 
type 2 is given by 
$y_{2t}(s^t)=1-y_{1t}(s^t)$.



\vskip.75cm

\begin{itemize}

\item[e.] \textbf{(30 points)} 
Compute a sequential trading equlibrium, i.e.,
find an allocation, prices 
$\{Q_t(s_{t+1}\vert s^t); \forall t, \forall s_{t+1}, \forall s^t\}$
and asset holdings  
$\{a_{it+1}(s_{t+1}, s^t); \forall i, \forall t+1, 
                          \forall s_{t+1}, \forall s^t\}$
expressed in terms of primitives. (Hint: Use  
your results above for a time $0$ trading equilibrium.) 
\medskip

\item[f.] \textbf{(5 points)} 
Use a no-arbitrage argument to find an
expression for the rate of return on a risk-free asset 
between periods
$t$ and $t+1$, $R_t(s^t)$. That is, express the sequence 
$\{R_t(s^t); \forall t, \forall s^t\}$ in terms of other 
prices.
(Provide a mathematical derivation rather than a verbal 
argument.)
%\medskip

%\item[g.] \textbf{(5 points)} 
%Suppose that the economy consists only
%of consumers of type~1. Compute a sequential trading 
%equilibrium for such an economy, and compare its prices, 
%including the risk-free 
%interest rate, to those of the original economy.

\end{itemize}

\vskip.5cm

\noindent
{\bf From hereon}, we drop the assumption of complete markets.
Instead, we assume that consumers can only borrow and
lend among themselves at a risk-free interest rate.
They trade in ``discount bonds'' as follows. 
Let $b_{it}(s^t)$ be the purchase of bonds at
time $t$ and history $s^t$ by consumers of type~$i$ 
at the price $p_t(s^t)$. 
Each unit of bonds promises to pay one consumption 
good next period $t+1$ regardless of which state 
$s_{t+1}$ is realized. A
negative value of $b_{it}(s^t)$ means that consumers 
of type~$i$ are issuing bonds and hence, they are
borrowing from others. The bonds are in zero net 
supply, $b_{1t}(s^t) + b_{2t}(s^t)=0$.

\vskip.75cm

\begin{itemize}

\item[g.] \textbf{(25 points)} 
Compute a sequential trading equlibrium in this
incomplete-markets economy, i.e.,
find an allocation, prices 
$\{p_t(s^t); \forall t, \forall s^t\}$
and asset holdings  
$\{b_{it}(s^t); \forall i,$ $\forall t, 
                          \forall s^t\}$
expressed in terms of primitives. (Hints: (1) Think
recursively. (2) Make the conjecture that 
$p_t(s^t)=\beta$
%for all $t$ and $s^t$, 
and later confirm that to be an equilibrium 
outcome.)

\medskip

\item[h.] \textbf{(5 points)} 
Compare consumption outcomes across the
two economies in questions e and g, including the
time $0$ expectations of future consumption
levels. 
Are there any agents who are better off in one
economy as compared to the other one? 

\end{itemize} 




\end{document}

